\chapter{Conclusion}\label{ch:conclusions}

\section{Key findings}

This thesis investigated \emph{how} visual instruction tuning alters the internal representations of LLaVA's LLM backbone and \emph{where} in depth those changes concentrate. We first showed that visual instruction tuning improves zero-shot QA performance, with the largest gains on open-ended, multimodal inputs. To localize where the backbone changes, we compared residual-stream representations between fine-tuned and pretrained checkpoints using Neighborhood Overlap and linear CKA at both the \emph{layer} and \emph{head} levels. Similarities drop most in early-mid layers---especially for multimodal inputs---with partial convergence across modalities in deeper layers for QA (Figures~\ref{fig:layer_wise_similarity_cocovqa}, \ref{fig:heads_wise_similarity_cocovqa}, \ref{fig:layer_wise_similarity_cocovqa_mean_heads}). 

We then examined geometric changes via intrinsic dimension and dataset entropy (Figures~\ref{fig:intrinsic_dimension}, \ref{fig:dataset_entropy}). Fine-tuning \emph{compresses} representations (lower ID/entropy) primarily in early-mid layers---again most strongly for multimodal inputs---while deep layers converge to similar, compact regimes across tasks. An analysis of the \emph{within-checkpoint} modality gap (Figures~\ref{fig:cosine_similarity_text_image}, \ref{fig:homogeneity_score_text_image}) showed slightly higher text-image alignment for the fine-tuned model in shallow layers, a small reversal in late layers, and task-dependent clustering behavior; taken together, these do not yield a single, unified pattern for the modality gap. Finally, layer patching validated the representational findings: cumulatively replacing early-mid pretrained layers with their fine-tuned counterparts recovers most performance gains on visual QA and image captioning, while patching late layers yields diminishing returns (Figures~\ref{fig:accuracy_question_answering_cumulative}, \ref{fig:cider_clip_score_image_captioning_cumulative}). 


Practically, these results suggest targeting early-mid layers during fine-tuning to capture multimodal integration, while keeping later layers stable to preserve next-token decoding. The patching outcomes also indicate that the multimodal connector learned during feature alignment may suffice for mapping the vision encoder features to the text embedding space of the language model, with most downstream task gains arising from editing the LLM backbone during fine-tuning. Interpretability efforts could also benefit from a focus on the early-mid layers to understand the multimodal integration.

\section{Limitations and Future Work}

This study focused on the impact of visual instruction tuning on the backbone representations of LLaVA, under the following constraints. First, we analyze only the LLaVA-1.5-7B model; the 13B version remains untested, and evaluating it would clarify whether our conclusions generalize. Second, our analysis is restricted to last-token residual-stream embeddings, a broadly used representation in the literature and consistent with our experimental design. Third, the results are limited to COCO-derived data for question answering and image captioning. Extending the analysis to other tasks and datasets would test the robustness of our findings. Finally, all experiments are based on the one-epoch fine-tuned LLaVA model released by the authors. The task-specific outcomes of our modality gap analysis (Figures~\ref{fig:accuracy_question_answering_cumulative} and \ref{fig:cider_clip_score_image_captioning_cumulative}) may reflect this constraint, and longer training schedules could yield different results. These limitations define the scope of the present study while pointing toward clear directions for future work on how visual instruction tuning reshapes multimodal LLMs and how to fine-tune them effectively and efficiently.